% Generated draft: requires mathematical and source review.
\documentclass[11pt]{article}
\usepackage[margin=25mm]{geometry}
\usepackage{fontspec}
\setmainfont{Latin Modern Roman}
\usepackage{amsmath,amssymb,mathtools,unicode-math}
\setmathfont{Latin Modern Math}
\usepackage{enumitem,needspace,xurl,hyperref}
\setlength{\parindent}{0pt}
\setlength{\parskip}{5pt}
\setlength{\emergencystretch}{3em}
\newcommand{\N}{\mathbb{N}}
\newcommand{\Z}{\mathbb{Z}}
\newcommand{\Q}{\mathbb{Q}}
\newcommand{\R}{\mathbb{R}}
\newcommand{\C}{\mathbb{C}}
\newcommand{\F}{\mathbb{F}}
\newcommand{\A}{\mathbb{A}}
\newcommand{\PP}{\mathbb{P}}
\newcommand{\E}{\mathbb{E}}
\newcommand{\eps}{\varepsilon}
\newcommand{\dd}{\,\mathrm{d}}
\begin{document}
\section*{UnsolvedMath ID: 701614}
\textbf{Source problem code:} catalogue entry
\subsection{Definitions and mathematical statement}
A variety of groups is a class of groups defined by identities in the language of groups. For a variety \(\Theta\) and a positive integer \(r\), let \(F_\Theta(r)\) denote the relatively free group in \(\Theta\) on \(r\) generators. Groups \(G\) and \(H\) are elementarily equivalent, written \(G\equiv H\), if they satisfy exactly the same first-order sentences in the language of groups.

The variety \(\Theta\) is of \emph{Tarski type} if
\[
\forall r,s\in\Z_{\geq 1},\qquad
\bigl(F_\Theta(r)\text{ and }F_\Theta(s)\text{ are non-abelian}\bigr)
\Longrightarrow F_\Theta(r)\equiv F_\Theta(s).
\]
The non-abelian restriction is part of the definition. The source does not require that any such non-abelian relatively free groups exist; whether examples with this additional non-vacuity condition are intended is unresolved.

For a positive integer \(n\), define the Burnside variety
\[
B_n=\{G: \forall x\in G,\ x^n=1\}.
\]
Here the identity requires every element order to divide \(n\); it does not require the least exponent of every group to equal \(n\). Using the convention \([a,b]=a^{-1}b^{-1}ab\), put
\[
e_1(x,y)=[x,y],\qquad e_{k+1}(x,y)=[e_k(x,y),y]\quad(k\geq 1).
\]
The \(n\)-Engel variety is
\[
E_n=\{G: \forall x,y\in G,\ e_n(x,y)=1\}.
\]
Thus the Engel word is left-normed, with \(n\) occurrences of \(y\). This makes explicit the conventional reading of the source's compressed displayed identity.

\begin{enumerate}
\item[(a)] Find examples of varieties \(\Theta\) of Tarski type distinct from the variety of all groups.
\item[(b)] Is \(B_n\) of Tarski type for sufficiently large \(n\)? Under the usual eventual interpretation, this asks whether
\[
\exists N_B\in\Z_{\geq 1}\ \forall n\in\Z_{\geq 1},\qquad
n\geq N_B\Longrightarrow B_n\text{ is of Tarski type}.
\]
\item[(c)] Is \(E_n\) of Tarski type for sufficiently large \(n\)? Under the same interpretation, this asks whether
\[
\exists N_E\in\Z_{\geq 1}\ \forall n\in\Z_{\geq 1},\qquad
n\geq N_E\Longrightarrow E_n\text{ is of Tarski type}.
\]
\end{enumerate}
In (b) and (c), the source says only that \(n\) is ``big enough.'' The displayed eventual quantifiers state a conventional interpretation, not a verified clarification from the proposers; no restriction to a special class of integers is supplied.
\subsection{Short English statement}
A variety is of Tarski type when all its non-abelian relatively free groups of finite rank have the same first-order theory. The problem asks for examples other than the variety of all groups, and asks whether the Burnside and Engel varieties have this property for sufficiently large parameters. The intended non-vacuity requirement and the precise scope of ``sufficiently large'' remain unclarified.
\subsection{Sources}
\begin{itemize}
\item Supplied local UnsolvedMath registry record, permanent ID \(701614\), exact source code \texttt{catalogue entry}, title \emph{Kourovka Notebook Problem 19.66}; upstream revision \texttt{d371ee917ac9647fbb06cd6fc529deec6f3dda82}. The record was created and updated on 2026-10-16.
\item The supplied record attributes Problem 19.66 to B. I. Plotkin and E. B. Plotkin, proposed in 2018. Its attribution field contains an unmatched LaTeX formatting command; the names are transcribed here without that command.
\item Supplied source attribution: \emph{The Kourovka Notebook (Unsolved Problems in Group Theory)}, 19th issue (2018), and 21st issue edition (2026), identified in the record as \texttt{arXiv:1401.0300v48}. The collection description names E. I. Khukhro and V. D. Mazurov as editors. These bibliographic details are supplied attribution and have not been independently checked.
\item The registry labels the problem \emph{open}. This is a supplied status label, not an independently verified literature assessment. No separate literature sources or literature-check date are supplied.
\end{itemize}
\end{document}
